\section{Integral path generators and Frobenius columns}
\label{sec:path-generators}

\status{Both generator assertions retained in [SY6] are proved below. The same integral filtered basis supplies the \texttt{SPAN} input to the local Gram estimate. The detailed companion proves the combinatorial and Frobenius calculations with the literal source conventions. No integral decomposition into rational isotypic summands is asserted.}

\subsection{One integral basis for all prefix conditions}

Give letters $x_0,\ldots,x_{N-1}$ in dimension $n$ vector or covector signs and degrees $d_i$. A signed Pieri path is a sequence of decreasing integer tuples $\lambda^0=0,\ldots,\lambda^N$: a vector step increases by a horizontal strip of size $d_i$, and a covector step decreases by a horizontal strip of size $d_i$. Set $h_i(c)=\#\{a:c\le\lambda^i_a\}$. Each column starts at $1_{\varnothing}$; when $c>0$, a height increase applies exterior multiplication by $x_i$, a decrease applies contraction. For $c\le0$ interchange those operations. Use the source's right-wedge signs. Call the resulting component $K_c(I_c)$. Choose $|I_c|=h_N(c)$ for $c>0$, and $|I_c|=n-h_N(c)$ otherwise. Complement the negative-column subsets. Require the increasing lists in these modified columns to be weakly increasing along rows. Put
\[
 \Psi_{\lambda,I}=\prod_cK_c(I_c).
 \tag{IP.1}
\]
Only finitely many factors differ from $1$. These are exactly the source's \texttt{Kcol}, \texttt{IsSt}, and \texttt{Psi}. \textsf{[IP1]}

\begin{lemma}[Integral filtered path basis]
The polynomials (IP.1) form a basis over $\mathbb Z$. Over any integral domain $R$, the simultaneous conditions
\[
 v_i(F)\ge\theta_i,\qquad M_i(F)\le U_i
 \tag{IP.2}
\]
are equivalent to an $R$-linear combination of the same basis elements satisfying $\lambda^i_n\ge\theta_i$ and $\lambda^i_1\le U_i$ for every prefix $i$. Here $v_i,M_i$ are the minimum and maximum weights of the generic conjugated rank-one torus on that prefix.
\end{lemma}
\begin{proof}
For canonical terminal subsets, order the last letter's exponents with coordinate $n$ first for vectors, coordinate $1$ first for covectors. The leading component is
\[
 \pm\Psi_{\text{preceding path, canonical}}
       \prod_a x_{N-1,a}^{|\lambda^N_a-\lambda^{N-1}_a|}.
\]
For fixed endpoint its exponent determines the preceding endpoint; induction proves canonical-path independence over every field. A generic matrix separates the other terminal choices: its endpoint-weight coefficient is the canonical product times a product of minors. Jacobi's complementary-minor identity turns this into a common determinant power times flag minors. Their diagonal leading monomials recover the weakly ordered tableau rows, so they are distinct. Taking a maximal endpoint in dominance proves independence of all (IP.1), in every characteristic. \textsf{[IP2]}

Row insertion of a weakly increasing word gives
\[
 \sum_{\mu/\lambda\text{ horizontal strip of size }d}s_n(\mu)
   =s_n(\lambda)\binom{n+d-1}{d},
\]
where $s_n$ counts standard terminal choices. Determinant shifts permit negative weights; complementing and reversing columns gives the downward version. Iteration counts exactly the ordinary monomials. The square integer coefficient matrix is thus nonsingular modulo every prime, so its determinant is $\pm1$. \textsf{[IP3]}

For the full prefix, every transformed column product has coordinate weights between $\lambda_n$ and $\lambda_1$. Conversely, a maximal bad endpoint in a basis expansion survives in its generic endpoint-weight component by the preceding independence argument. A good endpoint cannot dominate it. For an earlier prefix, cut each exterior recurrence there: this proves containment for good paths. Iterating the same counting identity from that endpoint gives exactly the dimension of the constrained prefix space times the tail monomial dimension. Thus containment is equality. Uniqueness of basis expansion handles all prefixes simultaneously. Passing to the fraction field and back retains integral coefficients. \textsf{[IP4]}
\end{proof}

Each column is multilinear exactly in its used letters, and letter $i$ occurs in $d_i$ columns. A positive column's height is at most its preceding vector occurrences. Cutting at prefix $i$ also gives
\[
 v_i(K_c(I_c))\ge q(c,h_i(c)),\quad
 \lambda^i_n=\sum_cq(c,h_i(c)),\quad
 q(c,h)=\mathbf1_{c>0,\,h=n}-\mathbf1_{c\le0,\,h<n}.
 \tag{IP.3}
\]
Indeed the intermediate exterior weights are $0,1$ for positive columns and $0,-1$ for negative ones. These facts are precisely the remaining \texttt{SPAN} clauses. \textsf{[IP5]}

\subsection{The actual generators and their keys}

For the source frame segment, an \texttt{MGen0} generator is $AG$ in blocks of dimensions $r,K$. Apply the lemma separately. The defining zero-gap inequalities give, for every surviving path pair $(\nu,\kappa)$,
\[
 \kappa^{i+1}_1\le\nu^i_r\quad(\text{vector}),\qquad
 \kappa^i_1\le\nu^{i+1}_r\quad(\text{covector}).
\]
On the initial preceding-frame block, $G$ has degree zero and $A$ uses only vectors. These same inequalities follow there. They are exactly the missing cross-block strip inequalities, so $(\nu,\kappa)$ is a genuine joined path. Multiplying integral expansions proves \texttt{GeneratorSpanInclusion}. \textsf{[IP6]}

For a joined path put $S_c=K_c^{\nu}(I'_c)K_c^{\kappa}(I''_c)$. The literal split-product identity gives $P=\prod_cS_c$. Let selected lifted positions be $s_0<\cdots<s_{N-1}\le z$, with $p^{z-s_0}\le Q=p^{G'}$. A step is costly when its height rises from $h\ge r$ or falls from $h\le r$. Set
\[
 a_c=\sum_{i\text{ costly}}p^{z-s_i}<2Q.
 \tag{IP.4}
\]
Every monomial of $S_c$ has this weighted normal degree. The source bound $N+r\le K$ keeps negative columns strictly above height $r$. \textsf{[IP7]}

Work over the source DVR $\mathcal O$, with $w^2=p$, $p\ge13$. Construct $F_c^g$ using the same exterior recurrence in the adapted chart-$g$ letters, but raise every component to $p^{s-t}$ between successive used positions, then to $p^{z-t}$ at the end. This is the integral polynomial \texttt{KKc}. Its degree at used position $s_i$ is $p^{z-s_i}$. \textsf{[IP8]}

Here is the essential all-chart calculation. In a moving integral exterior basis, write $VL_s=w^aZ$. Witt Frobenius gives, for $P=p^{t-s}$,
\[
 V^{[P]}L_t=w^{aP}(Z^{[P]}+pR),\qquad R\text{ integral}.
 \tag{IP.5}
\]
Hence an extra factor $w$ becomes two extra factors across a strict gap. The nominal subsets are first-block subsets below height $r$, subsets containing that block above it, and last-block subsets for negative columns. Inductively every component has the predicted order; nonnominal components have one extra. At a costly step, normal terms acquire $w$, while tangential terms require a nonnominal predecessor and retain two extra powers. This proves order $a_c$ in every chart. \textsf{[IP9]}

In its own chart, divide by the predicted power and reduce modulo $w$. Exactly the prescribed block operation survives, with its original exterior sign. Frobenius distributes in the residue field, giving
\[
 \Phi_gF_c^g\equiv w^{a_c}\operatorname{phSeg}(S_c)
                 \pmod{w^{a_c+1}}.
 \tag{IP.6}
\]
Thus each factor is a literal key, with the required selected-position support. For $o=\sum_ca_c$, their leading terms multiply to $w^o\operatorname{phSeg}(P)$, whose normal degree is $o$. This is exactly \texttt{Sy}, proving \texttt{GeneratorImageIsSymbol}. Zero factors cause no exception. \textsf{[IP10--IP11]}
